
\bibitem{src-1} M. Abouzaid, Nearby Lagrangians with vanishing Maslov class are homotopy equivalent, arXiv:1005.0358.

\bibitem{src-2} R. Abraham and J.W. Robbin, \emph{Transversal Mappings and Flows}, Benjamin, New York–Amsterdam, 1970.

\bibitem{src-3} M. Abreu, Topology of symplectomorphism groups of $S^{2}\times S^{2}$, \emph{Inv. Math.}, 131 (1998), 1–23.

\bibitem{src-4} M. Abreu, G. Granja, and N. Kitchloo, Compatible complex structures on symplectic rational ruled surfaces, \emph{Duke Math.J.} 148 (2009), no. 3, 539–600.

\bibitem{src-5} M. Abreu and D. McDuff, Topology of symplectomorphism groups of rational ruled surfaces, SG/9910057, \emph{J. Amer. Math. Soc.} 13, (2000) 971–1009.

\bibitem{src-6} M. Abreu and L. Macarini, Remarks on Lagrangian intersections in toric manifolds, arXiv:1105.0640.

\bibitem{src-7} L. Abrams, The quantum Euler class and the quantum cohomology of the Grassmannian, \emph{Israel Math.J.} 117 (2000), 335–352.

\bibitem{src-8} B. Aebischer, M. Borer, Ch. Leuenberger, M. Kälin, and H.M. Reimann, \emph{Symplectic geometry, an introduction based on the Seminar at Berne, 1992}, Birkhäuser, Basel, 1994.

\bibitem{src-9} S. Agnihotri, \emph{Quantum Cohomology and the Verlinde Algebra}, PhD thesis, Oxford University, September 1995.

\bibitem{src-10} A. Akhmedov and B.D. Park, Exotic smooth structures on small 4-manifolds with odd signatures, \emph{Invent. Math.} 181 (2010), 577–603.

\bibitem{src-11} M. Akveld and D.A. Salamon, Loops of Lagrangian submanifolds and pseudoholomorphic disks, \emph{Geom. Funct. Anal.} 11 (2001), 609–650.

\bibitem{src-12} P. Albers and H. Hofer, On the Weinstein conjecture in higher dimensions, \emph{Comm. Math. Helv.} 84 (2009), no 2, 429–436.

\bibitem{src-13} S. Anjos, The homotopy type of symplectomorphism groups of $S^{2}\times S^{2}$, \emph{Geom.\ and Top.} 6 (2002), 195–218.

\bibitem{src-14} S. Anjos and G. Granja, Homotopy decomposition of a group of symplectomorphisms of $S^{2}\times S^{2}$, \emph{Topology} 43 (2004) (3), 599–618.

\bibitem{src-15} V.I. Arnold, \emph{Mathematical methods in classical mechanics}, Springer, Berlin, 1978.

\bibitem{src-16} V.I. Arnold, Sur une proprieté topologique des applications globalement canoniques et à mécanique classique, C.R. Acad. Sci. Paris 261 (1965), 3719–3722.

\bibitem{src-17} V.I. Arnold, On a characteristic class entering into conditions of quantization, \emph{Functional analysis} 1 (1967), 1-8.

\bibitem{src-18} N. Aronszajn, A unique continuation theorem for elliptic differential equations or inequalities of the second order, \emph{J. Math. Pures Appl.} 36 (1957), 235–249.

\bibitem{src-19} P.S. Aspinwall and D.R. Morrison, Topological Field Theory and rational curves, hep-th/9110048, \emph{Comm. Math. Physics} 151 (1993), 245–262.

\bibitem{src-20} M.F. Atiyah, New invariants for three and four dimensional manifolds, \emph{Proc. Symp. Pure Math.} 48 (1988), 285–299.

\bibitem{src-21} M.F. Atiyah and I.G. MacDonald, \emph{Introduction to Commutative Algebra}, Addison Wesley, 1969.

\bibitem{src-22} M. Audin, Fibrés normaux d’immersions en dimension moitié, points doubles d’immersions lagrangiennes et plongements totalement réels, \emph{Comment. Math. Helv.} 63 (1988), 593-623.

\bibitem{src-23} M. Audin, \emph{The Topology of Torus Actions on Symplectic Manifolds}, Progress in Mathematics 93, Birkhäuser, 1991.

\bibitem{src-24} M. Audin and M. Damian, \emph{Théorie de Morse et homologie de Floer}, Savoirs Actuels (les Ulis), EDP Sciences, Les Ulis, CNRS Éditions, Paris 2010.

\bibitem{src-25} M. Audin and F. Lafontaine, ed., \emph{Holomorphic Curves in Symplectic Geometry}, Progress in Mathematics 117, Birkhäuser 1994.

\bibitem{src-26} D. Auroux, Symplectic 4-manifolds as branched coverings of $\mathbb{C}P^{2}$, \emph{Invent. Math.} 139 (2000), 551–602.

\bibitem{src-27} D. Auroux, Monodromy invariants in symplectic topology, SG/0304113, notes of lectures given at IPAM, Los Angeles.

\bibitem{src-28} D. Auroux and L. Katzarkov, Branched coverings of $\mathbb{C}P^{2}$ and invariants of symplectic 4-manifolds, \emph{Invent. Math.} 142 (2000), 631–673.

\bibitem{src-29} D. Auroux, V. Muñoz, and F. Presas, Lagrangian submanifolds and Lefschetz pencils, \emph{J. Symplectic Geom.} 3 (2005), no 2, 171–219.

\bibitem{src-30} S. Baldridge and P. Kirk, A symplectic manifold homeomorphic but not diffeomorphic to $\mathbb{C}P^{2}\#3\overline{\mathbb{C}P}^{2}$, \emph{Geom. Topol.} 12 (2008), 919–940.

\bibitem{src-31} V. Batyrev, Quantum cohomology rings of toric manifolds, \emph{Astérisque} 218 (1993), 9–34.

\bibitem{src-32} S. Bauer, Almost complex 4-manifolds with vanishing first Chern class, \emph{J. Differential Geom.} 79 (2008), (1), 25–32.

\bibitem{src-33} A. Beauville, \emph{Surfaces Algébriques Complexes}, \emph{Astérisque} 54, 1978.

\bibitem{src-34} K. Behrend and Yu. Manin, Stacks of stable maps and Gromov-Witten invariants, \emph{Duke Math. J.} 85 (1996), 1–60.

\bibitem{src-35} R. Berndt, \emph{An Introduction to Symplectic Geometry}, Graduate Studies in Math.\ vol 26, A.M.S. (2001).

\bibitem{src-36} A. Bertram, G. Daskalopoulos, and R. Wentworth, Gromov invariant for holomorphic maps from Riemann surfaces to Grassmannians, \emph{J. Amer. Math. Soc.} 9 (1996), 529–571.

\bibitem{src-37} M. Bialy and L. Polterovich, Geodesics of Hofer’s metric on the group of Hamiltonian diffeomorphisms, \emph{Duke Math. J.} 76 (1994), 273–292.

\bibitem{src-38} P. Biran, Symplectic packing in dimension 4, \emph{Geometric and Functional Analysis} 7 (1997), 420–437.

\bibitem{src-39} P. Biran, A stability property of symplectic packing, \emph{Invent. Math.} 136 (1999), 123–155.

\bibitem{src-40} P. Biran, From symplectic packing to algebraic geometry and back, European Congress of Mathematics, Barcelone 2000, 507–524, in Progress in Math. 202, Birkhäuser, Basel, 2001.

\bibitem{src-41} P. Biran and K. Cieliebak, Symplectic topology on subcritical manifolds, \emph{Comm. Math. Helv} 76 (2001), 712–753.

\bibitem{src-42} P. Biran and K. Cieliebak, Lagrangian embeddings into subcritical Stein manifolds, \emph{Israel J. Math.}, 127 (2002), 221-244.

\bibitem{src-43} P. Biran and O. Cornea, Rigidity and uniruling for Lagrangian submanifolds, \emph{Geom. Topol.} 13 (2009), no. 5, 2881–2989.

\bibitem{src-44} P. Biran and O. Cornea, A Lagrangian Quantum homology, New perspectives and challenges in symplectic field theory, 1–44, CRM Proc. Lecture Notes 49, Amer. Math. Soc., Providence, RI, 2009.

\bibitem{src-45} R. Bott and L. Tu, \emph{Differential Forms in Algebraic Topology}, Springer, New York, Heidelberg, Berlin 1982.

\bibitem{src-46} F. Bourgeois, Y. Eliashberg, H. Hofer, K. Wysocki and E. Zehnder, Compactness results in Symplectic Field theory, \emph{Geom.\ and Top.} 7 (2003), 799-833.

\bibitem{src-47} S. Bradlow, Special metrics and stability for holomorphic bundles with global sections, \emph{J. Diff. Geom.} 33 (1991), 169–213.

\bibitem{src-48} J. Bryan and R. Pandharipande, Curves in Calabi–Yau 3-folds and topological quantum field theories, \emph{Duke Math. J.} 126 (2005), no. 2, 369–396.

\bibitem{src-49} A. Buch, Quantum cohomology of partial flag manifolds, \emph{Trans. Amer. Math. Soc.} 357 (2005), no. 2, 443-458.

\bibitem{src-50} L. Buhovsky, The Maslov class of Lagrangian tori and quantum products in Floer cohomology, \emph{J. Topol. Anal.} 2 (2010), 57-75.

\bibitem{src-51} L. Buhovsky and Y. Ostrover, On the uniqueness of Hofer’s Geometry, arXiv:1008.06661.

\bibitem{src-52} L. Buhovsky, M. Entov and L. Polterovich, Poisson brackets and symplectic invariants, arXiv:1102.3198.

\bibitem{src-53} L. Buhovsky and S. Seyfaddini, Uniqueness of generating Hamiltonian for continuous Hamiltonian flows, arXiv:1003.2612.

\bibitem{src-54} O. Buse, Relative family Gromov–Witten invariants and symplectomorphisms, \emph{Pacific J. Math.} 218 (2005), no. 2, 315–341.

\bibitem{src-55} O. Buse, Negative inflation and stability in symplectomorphism groups of ruled surfaces, \emph{J. Symplectic Geom.} 9 (2011), no. 2, 147–160.

\bibitem{src-56} O. Buse and R. Hind, Symplectic embeddings of ellipsoids in dimension greater than four, arXiv:1112.1452.

\bibitem{src-57} P. Candelas and X.C. de la Ossa, Moduli space of Calabi-Yau manifolds, \emph{Nuclear Physics} B355 (1991), 455–481.

\bibitem{src-58} P. Candelas, X.C. de la Ossa, P.S. Green, and L. Parkes, A pair of Calabi-Yau manifolds as an exactly soluble superconformal field theory, \emph{Nuclear Phys B} 359 (1991), 21–74.

\bibitem{src-59} A. Cannas da Silva, \emph{Lectures on Symplectic Geometry}, Lecture Notes in Mathematics vol 1764, Springer, Berlin, 2001.

\bibitem{src-60} P. Cascini and D. Panov, Symplectic generic complex structures on 4-manifolds with $b^{+}=1$, arXiv:1012.3644.

\bibitem{src-61} Yu. Chekanov, Invariant Finsler metrics on the space of Lagrangian embeddings, \emph{Math. Zeitschrift} 234 (2000), 605–619.

\bibitem{src-62} W. Chen and Y. Ruan, Orbifold Gromov–Witten theory, AG/0103156, in: Orbifolds in Mathematics and Physics (Madison 2001), Contemp. Math. 310 Amer. Math. Soc. (2002), 25–85.

\bibitem{src-63} K. Cieliebak and Ya. Eliashberg, \emph{From Stein to Weinstein and back: Symplectic geometry of affine complex manifolds}, Amer. Math. Soc., Providence, RI, 2012.

\bibitem{src-64} K. Cieliebak, A.R. Gaio, D.A. Salamon, $J$-holomorphic curves, moment maps, and invariants of Hamiltonian group actions, \emph{Int. Math. Res. Notes} 10 (2000), 831–882.

\bibitem{src-65} K. Cieliebak, H. Hofer, J. Latschev and F. Schlenk, Quantitative symplectic geometry, Dynamics, Ergodic Theory, Geometry MSRI, 54 (2007), 1–44.

\bibitem{src-66} K. Cieliebak, I. Mundet i Riera and D.A. Salamon, Equivariant moduli problems and the Euler class, \emph{Topology} 42 (2003), 641–700.

\bibitem{src-67} K. Cieliebak, R.A. Gaio, I. Mundet i Riera and D.A. Salamon, The symplectic vortex equations and invariants of Hamiltonian group actions, \emph{J. Sympl. Geom.} 1 (2002), 543–645.

\bibitem{src-68} K. Cieliebak and K. Mohnke, Symplectic hypersurfaces and transversality in Gromov–Witten theory, \emph{J. Symplectic Geom.} 3 (2005), no. 4, 589–654.

\bibitem{src-69} K. Cieliebak and D.A. Salamon, Wall crossing for symplectic vortices and quantum cohomology, \emph{Math. Ann.} 335 (2006), no. 1, 133–192.

\bibitem{src-70} I. Ciocan-Fontanine, Quantum cohomology of flag varieties, \emph{Int. Math. Res. Notes}, 2 (1995), 263–277.

\bibitem{src-71} T. Coates and A.B. Givental, Quantum Riemann–Roch, Lefschetz and Serre, \emph{Ann.\ of Math.} (2) 165 (2007), no. 1, 15–53.

\bibitem{src-72} T. Coates, Riemann-Roch theorems in Gromov–Witten theory, Ph.D. thesis, Berkeley (2003).

\bibitem{src-73} R. Cohen, J.D.S. Jones and G. Segal, Floer’s infinite dimensional Morse theory and homotopy theory, in [181], 297–325.

\bibitem{src-74} V. Colin, P. Ghiggini, and Ko Honda, Equivalence of Heegaard Floer homology and embedded contact homology via open book decompositions, \emph{Proc. Natl. Acad. Sci. USA} 108 (2011), no, 20, 8100–8105.

\bibitem{src-75} C. Conley and E. Zehnder, The Birkhoff–Lewis fixed point theorem and a conjecture of V.I. Arnold, \emph{Invent. Math.}, 73 (1983), 33–49.

\bibitem{src-76} D. Cox and S. Katz, \emph{Mirror Symmetry and Algebraic Geometry}, Math. Surveys and Monographs vol 68, AMS, Providence 1999.

\bibitem{src-77} M. Damian, Floer homology on the universal cover, a proof of Audin’s conjecture and other constraints on Lagrangian submanifolds. Preprint 2010, to appear in Comment. Math. Helv.

\bibitem{src-78} M. W. Davis and T. Januszkiewicz, Convex polytopes, Coxeter orbifolds and torus actions, \emph{Duke Math Journ.} 62 (1991), 417–451.

\bibitem{src-79} T. Delzant, Hamiltoniens périodiques et images convexes de l’application moment, \emph{Bull. Soc. Math. France} 116 (1988), 315–339.

\bibitem{src-80} R. Dijkgraaf, C. Faber, G. van der Geer, eds., \emph{The moduli space of curves}, Progress in Mathematics 129 Birkhäuser, 1995.

\bibitem{src-81} S.K. Donaldson, The orientation of Yang-Mills moduli spaces and 4-manifold topology, \emph{J. Diff. Geom.} 26 (1987), 397–428.

\bibitem{src-82} S.K. Donaldson, Gluing techniques in the cohomology of moduli spaces, in: \emph{Topological methods in modern mathematics (Stony Brook 1991)}, 137–170, Publish or Perish, TX, 1993.

\bibitem{src-83} S.K. Donaldson, Floer homology and algebraic geometry, in \emph{Vector Bundles in Algebraic Geometry (Durham 1993)}, 119–138, London Math. Soc. Lecture Notes 208, Cambridge Univ. Press, Cambridge 1995.

\bibitem{src-84} S.K. Donaldson, Symplectic submanifolds and almost-complex geometry, \emph{J. Diff. Geom.} 44 (1996), 666–705.

\bibitem{src-85} S.K. Donaldson, Lefschetz pencils on symplectic manifolds, \emph{J. Diff. Geom.} 53 (1999), 205–236.

\bibitem{src-86} S.K. Donaldson, Moment maps and diffeomorphisms, \emph{Asian J. Math.} 3 (1999), 1–15.

\bibitem{src-87} S.K. Donaldson, \emph{Floer homology groups in Yang–Mills Theory}, Cambridge University Press, Cambridge, 2002.

\bibitem{src-88} S.K. Donaldson and P. Kronheimer, \emph{The Geometry of Four-Manifolds}, Oxford University Press, Oxford, 1990.

\bibitem{src-89} S.K. Donaldson and I. Smith, Lefschetz pencils and the canonical class for symplectic 4-manifolds, \emph{Topology} 42 (2003), 743–785.

\bibitem{src-90} S. Dostoglou and D.A. Salamon, Instanton homology and symplectic fixed points, in \emph{Symplectic Geometry}, edited by D. Salamon, 57–93, LMS Lecture Notes Series 192, Cambridge University Press, 1993.

\bibitem{src-91} S. Dostoglou and D.A. Salamon, Cauchy–Riemann operators, self-duality, and the spectral flow, in \emph{First European Congress of Mathematics, Volume I, Invited Lectures (Part 1)}, edited by Joseph et al., 511–545, Birkhäuser, Progress in Mathematics, Vol. 119 (1994).

\bibitem{src-92} S. Dostoglou and D.A. Salamon, Self-dual instantons and holomorphic curves, \emph{Annals of Math.} 139 (1994), 581–640.

\bibitem{src-93} B. Dubrovin, Integrable systems in topological field theory, \emph{Nucl. Phys. B} 379 (1992), 627–689.

\bibitem{src-94} B. Dubrovin, Geometry of 2D topological field theories, in \emph{Springer Lecture Notes in Mathematics} 1620 (1996), 120–348.

\bibitem{src-95} C. Earle and J. Eells: A fibre bundle approach to Teichmüller theory, \emph{J. Diff. Geo.} 3 (1969), 19–43.

\bibitem{src-96} D. Eisenbud and W. Neumann, \emph{Three Dimensional Link Theory and Invariants of Plane Curve Singularities}, Annals of Mathematics Studies \# 110, Princeton Univ. Press, 1968.

\bibitem{src-97} I. Ekeland and H. Hofer, Symplectic topology and Hamiltonian dynamics, \emph{Math. Zeit.} 200 (1989), 355–378.

\bibitem{src-98} Y. Eliashberg, A theorem on the structure of wave fronts and its applications in symplectic topology, \emph{Functional Analysis and Applications} 21 (1987), 65–72.

\bibitem{src-99} Y. Eliashberg, D. Fuchs, T. Ratiu and A Weinstein, eds., \emph{Northern California Symplectic Geometry Seminar}, Amer. Math. Soc. Transl. (2) 196 (1999).

\bibitem{src-100} Y. Eliashberg and M. Gromov, Convex symplectic manifolds, in \emph{Several Complex Variables and Complex Geometry}, Part 2 (Santa Cruz CA 1989), Proc. Sympos. Pure Math. 52, Part 2, Amer. Math. Soc., Providence, RI (1991), 135–162.

\bibitem{src-101} Y. Eliashberg, A.B. Givental, and H. Hofer, Introduction to symplectic field theory, \emph{Geom. Funct. Anal.} 10 (2000), 560–673.

\bibitem{src-102} Y. Eliashberg and N. Mishachev, \emph{An introduction to the h-principle}, Graduate Texts \# 48, Amer. Math. Soc., Providence, RI, 2002.

\bibitem{src-103} Y. Eliashberg and L. Traynor, eds., \emph{Symplectic Geometry and Topology}, IAS/Park City Mathematics Series, Vol. 7, Amer. Math. Soc., Providence, RI, 1999.

\bibitem{src-104} M. Entov, $K$-area, Hofer metric and geometry of conjugacy classes in Lie groups, \emph{Invent. Math.} 146 (2001), 93–141.

\bibitem{src-105} M. Entov and L. Polterovich, Calabi quasimorphism and quantum homology, \emph{Int. Math. Res. Notes} (2003), no. 30, 1635–1676.

\bibitem{src-106} M. Entov and L. Polterovich, Rigid subsets of symplectic manifolds, \emph{Compos. Math.} 145 (2009), no. 3, 773–826.

\bibitem{src-107} M. Entov, L. Polterovich and F. Zapolsky, An “anit-Gleason” phenomenon and simultaneous measurement in Classical Mechanics, \emph{Foundations of Physics}, 37 (8), (1007)\sourcefootnote{原书题名作 “anit-Gleason”，年份作 “1007”；两者均疑为排印错误，但仅凭本书未能核定书目原题与准确年份，故照录。}, 1306–1316.

\bibitem{src-108} J.D. Evans, Symplectic mapping class groups of some Stein and rational surfaces, \emph{J. Symplectic Geom.} 9 (2011), no.1, 45–82.

\bibitem{src-109} J. Fine and D. Panov, The diversity of Calabi–Yau six-manifolds, arXiv:1108.5944.

\bibitem{src-110} R. Fintushel and R. Stern, Knots, links and 4-manifolds, \emph{Invent. Math.} 134 (1998), 363–400.

\bibitem{src-111} R. Fintushel and R. Stern, Pinwheels and nullhomologous surgery on 4-manifolds with $b^{+}=1$, \emph{Algeb. Geom. Topol.} 11 (2011), 1649–1699.

\bibitem{src-112} A. Floer, Proof of the Arnold conjecture for surfaces and generalizations to certain Kähler manifolds, \emph{Duke Math. J.} 53 (1986), 1–32.

\bibitem{src-113} A. Floer, Morse theory for Lagrangian intersections, \emph{J. Diff. Geom.} 28 (1988), 513–547.

\bibitem{src-114} A. Floer, The unregularized gradient flow of the symplectic action, \emph{Comm. Pure Appl. Math.} 41 (1988), 775–813.

\bibitem{src-115} A. Floer, Wittens complex and infinite dimensional Morse theory, \emph{J. Diff. Geom.} 30 (1989), 207–221.

\bibitem{src-116} A. Floer, Symplectic fixed points and holomorphic spheres, \emph{Comm. Math. Phys.} 120 (1989), 575–611.

\bibitem{src-117} A. Floer, An instanton invariant of 3-manifolds, \emph{Comm. Math. Phys.} 118 (1989), 215–240.

\bibitem{src-118} A. Floer and H. Hofer, Coherent orientations for periodic orbit problems in symplectic geometry, \emph{Math. Zeit.} 212 (1993), 13–38.

\bibitem{src-119} A. Floer, H. Hofer, and D.A. Salamon, Transversality in elliptic Morse theory for the symplectic action, \emph{Duke Math. J.} 80 (1995), 251–292.

\bibitem{src-120} U. Frauenfelder, Gromov convergence of pseudoholomorphic discs, Diploma thesis, ETH Zürich, February 2000.

\bibitem{src-121} U. Frauenfelder, Floer homology of symplectic quotients and the Arnold–Givental conjecture, PhD thesis, ETH-Zürich, January 2003.

\bibitem{src-122} U. Frauenfelder, The Arnold-Givental conjecture and moment Floer homology, \emph{IMRN} (2004), 2179–2269.

\bibitem{src-123} U. Frauenfelder and F. Schlenk, Hamiltonian dynamics on convex symplectic manifolds, \emph{Israel J. Math.} 159 (2007), 1–56.

\bibitem{src-124} A. Friedman, \emph{Partial Differential Equations}, Holt, Rinehart, and Winston, 1969.

\bibitem{src-125} K. Fukaya, Floer homology, $A^{\infty}$-categories, and topological field theory, in \emph{Proceedings of GARC Workshop on Geometry and Topology ‘93 (Seoul)}, ed H. Kim, Lecture Notes Ser 18, Seoul Nat. Univ, Seoul 1993, 1–102.

\bibitem{src-126} K. Fukaya, Floer homology for 3-manifolds with boundary, \emph{Topology, geometry and field theory}, 1–21, World Sci. Publ., River Edge, NJ, 1994. with update 1997 at \url{http://www.kusm.kyoto-u.ac.jp/fukaya/fukaya.html}.

\bibitem{src-127} K. Fukaya and K. Ono, Arnold conjecture and Gromov–Witten invariants for general symplectic manifolds, \emph{Topology} 38 (1999), 933–1048.

\bibitem{src-128} K. Fukaya, Y.-G. Oh, H. Ohta, K. Ono, \emph{Lagrangian Intersection Theory, Anomaly and Obstruction, Parts I and II}, AMS/IP Studies in Advanced Mathematics, Amer. Math. Soc. (Providence RI) and Internat. Press (Somerville MA), 2009.

\bibitem{src-129} K. Fukaya, Y.-G. Oh, H. Ohta, K. Ono, Lagrangian Floer theory on compact toric manifolds, I, \emph{Duke Math. J.} 151 (2010), no 1, 23-174.

\bibitem{src-130} K. Fukaya, Y.-G. Oh, H. Ohta, K. Ono, Lagrangian Floer theory on compact toric manifolds, II: bulk deformations, \emph{Selecta Math. (N.S.)} 17 (2011), no.3, 609–711.

\bibitem{src-131} K. Fukaya, P. Seidel and I. Smith, The symplectic geometry of cotangent bundles from a categorical viewpoint, in \emph{Homological mirror symmetry}, 1–26, Lecture Notes in Physics, vol 757, Springer, Berlin, 2009.

\bibitem{src-132} W. Fulton and R. MacPherson, A compactification of configuration space, \emph{Annals of Math.} 139 (1994), 183–225.

\bibitem{src-133} W. Fulton and R. Pandharipande, \emph{Notes on Stable Maps and Quantum Cohomology}, in \emph{Algebraic Geometry – Santa Cruz 1995}, eds. Kollar, Lazarsfeld and Morrison, Proc. Symp. Pure. Math. 62 part 2 (1997), 45–96.

\bibitem{src-134} A.R. Gaio, $J$-holomorphic curves and moment maps, PhD-thesis, University of Warwick, April 2000.

\bibitem{src-135} A.R. Gaio and D.A. Salamon, Gromov–Witten invariants of symplectic quotients and adiabatic limits, \emph{J. Symplectic Geom.} 3 (2005), no 1, 55–159. Erratum: J. Symplectic Geom. 7 (2009), no 3, 377–379.

\bibitem{src-136} O. Garcia-Prada, A direct existence proof for the vortex equations over a compact Riemann surface, \emph{Bull. London Math. Soc.} 26 (1994), 88–96.

\bibitem{src-137} R. Gautschi, Floer homology and surface diffeomorphisms, PhD thesis, ETH-Zürich, April 2002.
